\paragraph{Additional Published-Metric Comparisons.}
We also ran two lightweight, no-generation comparisons inspired by recent long-form and benchmarking work. First, following the SAGED emphasis on disparity diagnostics, we summarized maximum group disparity and a max group z-score over our existing audit metrics. For the primary \texttt{llm\_quality} metric, the SAGED-inspired summary matches the main result: baseline had a maximum race gap of 0.110 and max absolute group z-score of 0.153, while persona-blind reduced these to 0.017 and 0.069. Reranked and system-prompt conditions reduced the gap to 0.063 and 0.072, respectively. This comparison is not a reproduction of SAGED's full calibrated pipeline, but it applies a similar disparity-diagnostic framing to DemoGen's saved outputs.

Second, inspired by cover-letter bias transmission work such as Morehouse's GPT-4 cover-letter study, we checked whether cover-letter \texttt{llm\_quality} differed by perceived gender in our data. The gender gaps were small and non-significant across all conditions: baseline female-vs-male gap 0.022 ($p=0.757$), persona-blind gap 0.034 ($p=0.550$), reranked gap 0.007 ($p=0.928$), and system prompt gap 0.021 ($p=0.825$). This supports a narrower interpretation of DemoGen: the strongest evidence in our audit is race/name-linked quality disparity, not a broad claim that every demographic dimension or every task shows significant bias.
